% ===================================================================
\section{Reduction and Confluence}
\label{ch:reduction}
% ===================================================================

\subsection{Overview}

This chapter defines the syntactic layer of FRFP --- the Task Decomposition Grammar
(TDG) that generates pipeline terms --- and the semantic reduction rules that
normalize them. We prove two structural properties: \emph{termination} (every
reduction sequence is finite) and \emph{confluence} (the order of reductions does
not affect the final result). Together these establish that every pipeline term
has a unique normal form modulo navigation equivalence.

% -------------------------------------------------------------------
\subsection{Task Decomposition Grammar}
% -------------------------------------------------------------------

\begin{definition}[TDG terms]
A TDG term over the kernel primitives is a well-typed finite sequence:
\[
  t ::= \varepsilon \mid p \cdot t
\]
where $p$ is a primitive and $t$ is a TDG term, subject to the typing
constraint that consecutive primitives compose.

TDG terms are syntactic objects. A TDG term \emph{denotes} a pipeline (morphism)
in $\Czero$ when interpreted by the semantic function $\llbracket \cdot \rrbracket$
(Chapter~\ref{ch:semantics}).
\end{definition}

\begin{definition}[Free category]
The free category $\mathbf{F}$ generated by the six primitives has:
\begin{itemize}
  \item Objects: $\{\varnothing, \Eobj, \Tobj\}$
  \item Morphisms: all well-typed finite compositions of primitives
\end{itemize}
$\mathbf{F}$ is the syntactic model; $\Czero$ is its semantic quotient.
\end{definition}

% -------------------------------------------------------------------
\subsection{Reduction Relations}
% -------------------------------------------------------------------

\begin{definition}[Explicit reduction]
A reduction relation $\to_e$ on TDG terms captures explicit-space
simplifications:
\[
  \EC \cdot \EC \to_e \EC,\qquad
  \ED \cdot \EC \to_e \ED
\]
together with similar rules for all combinations of
$\{\RI, \EC, \ED\}$.
\end{definition}

\begin{definition}[Tacit reduction]
A reduction relation $\to_t$ on TDG terms captures tacit-space simplifications:
\[
  \TE \cdot \TE \to_t \TE,\qquad
  \HFD \cdot \TE \to_t \HFD
\]
\end{definition}

\begin{definition}[Combined reduction]
The combined reduction $\to_c = \to_e \cup \to_t$, with contextual closure.
The reflexive-transitive closure is written $\twoheadrightarrow_c$.
\leanref{Frfp.Core.OperationalSemantics}, \leanref{Frfp.Core.SemanticCorrectness}
\end{definition}

% -------------------------------------------------------------------
\subsection{Normal Forms}
% -------------------------------------------------------------------

\begin{definition}[Normal form]
A TDG term $t$ is in \emph{normal form} if no reduction rule applies. A normal
form has the structure:
\[
  t_{\mathrm{nf}} =
  \underbrace{e_1 \cdots e_k}_{\text{explicit block}} \cdot
  \underbrace{\RB^?}_{\text{optional boundary}} \cdot
  \underbrace{h_1 \cdots h_m}_{\text{tacit block}}
\]
where each $e_i \in \{\RI, \EC, \ED\}$, $\RB^?$ is either $\RB$ or absent,
and each $h_j \in \{\TE, \HFD\}$.

\leanref{Frfp.Minimal.BasisSweep.PrimitiveTraceAdmissibleStrong}
\end{definition}

This three-block structure is the \textbf{protocol normal form}. It reflects the
operational semantics of an FRFP interaction: explicit work first, then
(optionally) handoff to human tacit evaluation.

% -------------------------------------------------------------------
\subsection{Termination}
% -------------------------------------------------------------------

\begin{theorem}[Termination]
Every reduction sequence in $\to_c$ is finite.
\leanref{Frfp.Core.ConfluenceProof --- termination component}
\end{theorem}
\begin{proof}
Define a measure $\mu(t)$ as the number of consecutive same-phase pairs in $t$.
Each reduction step strictly decreases $\mu(t)$. Since $\mu(t) \geq 0$, all
sequences terminate. \qed
\end{proof}

% -------------------------------------------------------------------
\subsection{Local Confluence}
% -------------------------------------------------------------------

\begin{definition}[Local confluence]
The relation $\to_c$ is \emph{locally confluent} if for any $t, t_1, t_2$ with
$t \to_c t_1$ and $t \to_c t_2$, there exists $t_3$ with
$t_1 \twoheadrightarrow_c t_3$ and $t_2 \twoheadrightarrow_c t_3$.
\end{definition}

\begin{theorem}[Local confluence of $\to_c$]
The combined reduction is locally confluent.
\leanref{Frfp.Core.ConfluenceProof.diamond\_property}
\end{theorem}
\begin{proof}
By case analysis on pairs of one-step reductions from $t$. Since explicit and
tacit reductions act on disjoint parts of the term (by the ETS theorem), any two
one-step reductions are either in the same phase at non-overlapping positions
(trivially confluent), in the same phase at overlapping positions (resolved by the
phase-specific confluence lemma), or in different phases (independently applicable
in either order). \qed
\end{proof}

\begin{note}
Lifted action commutativity --- navigation actions commute with explicit reduction
--- is axiomatized from the operational semantics.
\leanref{Frfp.Core.Confluence.lifted\_action\_commutes}

This property is implicit in the original published theory. The Lean formalization
makes it an explicit required lemma.
\end{note}

% -------------------------------------------------------------------
\subsection{Global Confluence (Church--Rosser)}
% -------------------------------------------------------------------

\begin{theorem}[Confluence]
By Newman's Lemma, termination together with local confluence implies global
confluence:
\[
  \forall t, t_1, t_2,\quad
  t \twoheadrightarrow_c t_1 \wedge t \twoheadrightarrow_c t_2
  \Rightarrow
  \exists t_3,\quad
  t_1 \twoheadrightarrow_c t_3 \wedge t_2 \twoheadrightarrow_c t_3
\]
\leanref{Frfp.Core.ConfluenceProof.confluence}\\
\emph{Source}: Newman's Lemma \textbf{[Newman 1942]}; abstract reduction systems
\textbf{[Baader--Nipkow 1998]} Ch.~2.
\end{theorem}

% -------------------------------------------------------------------
\subsection{Normal Form Uniqueness (Corrected)}
% -------------------------------------------------------------------

\begin{theorem}[Normal form uniqueness modulo navigation equivalence]
For any TDG term $t$, all normal forms reachable from $t$ are navigation-equivalent:
\[
  \forall t,\ \forall \mathrm{nf}_1, \mathrm{nf}_2,\quad
  t \twoheadrightarrow_c \mathrm{nf}_1
  \wedge t \twoheadrightarrow_c \mathrm{nf}_2
  \wedge \mathrm{IsNF}(\mathrm{nf}_1)
  \wedge \mathrm{IsNF}(\mathrm{nf}_2)
  \Rightarrow
  \mathrm{nf}_1 \navequiv \mathrm{nf}_2
\]
\leanref{Frfp.Core.SemanticCorrectness.unique\_nf\_mod\_nav}
\end{theorem}

\begin{correction}[Lean Correction]
The original published theory states normal forms are \emph{literally unique}
(equal). The Lean formalization shows they are unique only \emph{modulo navigation
equivalence} $\navequiv$. Two normal forms from the same term may differ in
navigation annotations while denoting the same pipeline.
\end{correction}

\begin{definition}[Navigation equivalence]
Two terms $t_1 \navequiv t_2$ if they differ only in navigation morphisms
(admissible path annotations added by the navigation layer). Navigation
morphisms do not change the computational content of a term.
\leanref{Frfp.Core.SemanticCorrectness.NavEquivalent}
\end{definition}

% -------------------------------------------------------------------
\subsection{Semantic Correctness}
% -------------------------------------------------------------------

\begin{theorem}[Observable invariance]
For any term $t$ and its normal form $\mathrm{nf}$, the observable correctness
judgment is preserved:
\[
  \obs(t) = \obs(\mathrm{nf})
\]
where $\obs = \gamma \circ \Sem \circ \mathrm{initial}$ is the composition of
semantic evaluation, correctness judgment extraction, and trajectory initialization.
\leanref{Frfp.Core.SemanticCorrectness.obsCorrect\_invariant}
\end{theorem}

\begin{definition}[Correctness judgment type]
The correctness judgment type has three values:
\[
  \mathrm{CorrectnessJudgment} =
  \{\mathrm{correct},\ \mathrm{incorrect},\ \mathrm{unknown}\}
\]
This is a new type made explicit by the Lean formalization; it was implicit in
the original paper.
\leanref{Frfp.Core.TacitDependence.CorrectnessJudgment}
\end{definition}

% -------------------------------------------------------------------
\subsection{Protocol Normal Form Theorem}
% -------------------------------------------------------------------

\begin{theorem}[Protocol normal form]
Any strongly admissible primitive trace decomposes as an explicit block,
optional boundary crossing, and tacit block:
\[
  \forall \mathcal{T},\
  \mathrm{AdmissibleStrong}(\mathcal{T})
  \Rightarrow
  \exists e_1\ldots e_k,\ h_1\ldots h_m,\ b \in \{\top, \bot\},
\]
\[
  \mathcal{T}
  = [e_1, \ldots, e_k]
    \mathbin{+\!\!+}
    (\text{if } b \text{ then } [\RB] \text{ else } [])
    \mathbin{+\!\!+}
    [h_1, \ldots, h_m]
\]
\leanref{Frfp.Minimal.BasisSweep.protocol\_normal\_form\_theorem} ---
Machine-checked proof.
\end{theorem}
